\chapter{第十讲\quad 排除犯罪事由II：排除责任事由}
\label{lecture:10}
\noindent{\small\hyperref[ref:lecture:10]{查阅本讲扩展规则}}\par

先确认行为对应的违法事实，再评价行为人在行为时能否承担责任。年龄是按行为与罪行开启的门，责任能力评价辨认和控制。后两节讨论可避免的法律错误与行为时的合规选择空间，须区分法条和学理。

\section{第一节\quad 责任年龄}

\paragraph{法条和模型}
《刑法》第17条、第17条之一。设 \(a=\mathrm{Age}(p,t)\) 为行为人 \(p\) 实施行为 \(b\) 时已满周岁数，\(B_8(p,b)\) 表示其行为落入第17条第2款的八类罪行。\(I(p,b)\) 为故意杀人或故意伤害，\(D(p,b)\) 为可归责地致人死亡，\(K(p,b)\) 为特别残忍手段致重伤并造成严重残疾，\(E(p,b)\) 为该参与人情节恶劣，\(A(p,b)\) 为最高人民检察院核准追诉该人所涉行为。八类为故意杀人、故意伤害致人重伤或者死亡、强奸、抢劫、贩卖毒品、放火、爆炸、投放危险物质。年龄入口 \(Y\) 为
\[
Y(p,b,t)=
\begin{cases}
0, &a<12,\\
I(p,b)\land[D(p,b)\lor K(p,b)]\land E(p,b)\land A(p,b),
   &12\le a<14,\\
B_8(p,b), &14\le a<16,\\
1, &a\ge16.
\end{cases}
\]
这是在行为事实与核准状态已确定时的规则，证据不明须保留待核状态。\(Y=1\) 只确认年龄门，仍须证明构成与责任的其他条件。根据未成年人刑事案件司法解释，按公历计算并从周岁生日次日起算，故日期必须与行为点对齐。2026年高检发办字〔2026〕153号规定细化核准程序与恶劣情节判断。未经最高人民检察院核准，不得提起公诉。共同案件中的恶劣情节要按各参与人的地位和作用分别审查。因年龄不予刑事处罚的，依法另审管教与专门矫治教育。

\paragraph{推导}
设甲已满13周岁，实施故意伤害，特别残忍手段致乙重伤并造成严重残疾，且经综合评估情节恶劣。此时有 \(I\land K\land E\)，故第2行化为 \(Y=A\)。省级检察院同意报请时，不能推出 \(A\)。只有给定最高人民检察院核准，年龄入口才通过。若同案中的丙只是受诱导且作用、动机不同，甲的 \(E\) 不能直接赋给丙。若乙重伤但未严重残疾且未死亡，\(D\lor K=0\)，即使恶劣也不能满足该行。

\paragraph{例外和处罚}
15周岁人实施绑架并故意杀人，按罪行逐项分析，不能让绑架标签掩盖故意杀人。14至16周岁的人盗窃等后当场施暴，要查未成年人司法解释对事后转化抢劫的特殊规定，分别审故意重伤或故意杀人，不能直接把成人的转化规则套入 \(B_8\)。教材在部分帮助行为、轻伤故意与重伤结果的边界有专门分析，答题要先写行为、主观内容、结果，再判断是否落入八类。已经依法追责且未满18周岁，应当从轻或者减轻处罚。满75周岁，故意犯罪可以从轻或者减轻，过失犯罪应当从轻或者减轻。两项是量刑规则。

\section{第二节\quad 责任能力}

\paragraph{法条和模型}
《刑法》第18条、第19条。令 \(R(p,t)\) 为能辨认，\(C(p,t)\) 为能控制，\(F(p,t)\) 为精神障碍及法定鉴定等法定条件成立。对精神病人的完全丧失情形，有
\[
F(p,t)\land[\neg R(p,t)\lor\neg C(p,t)]
\Longrightarrow \text{该行为不负刑事责任}.
\]
限制责任能力须与完全丧失相区分，尚未完全丧失仍负刑事责任，可以从轻或者减轻。间歇性精神病人在精神正常时实施行为，按正常行为时状态评价。第18条关于醉酒犯罪仍负刑事责任，第19条关于又聋又哑或者盲人可以从轻、减轻或者免除处罚，均应在相应法定事实和构成要件下适用。

\paragraph{推导}
甲在精神正常时故意投毒，之后突发经鉴定确认的精神病，乙在发作期间中毒死亡。令 \(t_b\) 为投毒时、\(t_r\) 为死亡时，\(G\) 为投毒造成可归责死亡。由于
\[
R(\text{甲},t_b)\land C(\text{甲},t_b)\land G
\]
均已证成，故在其他要件齐备时可认定故意杀人既遂。把 \(t_r\) 的无能力反代 \(t_b\) 会偷换评价时点。若有能力时的一击没有造成可归责死亡，死亡独由失能后的另一次行为造成，有能力时的故意杀人仅在符合着手与未得逞条件时认定未遂。

\paragraph{两时点边界}
教材所说的原因自由行为，需分别证明先前可归责地造成失能及其与后来侵害的联系。例中若甲为伤害乙故意使自己陷入有充分证据证明的失能状态，后来实施计划内伤害，教材依原因自由行为讨论责任，并把结果行为的开始作为着手。该处理属于学理。若实际实行的是原计划未涵盖的抢劫，不得凭先前伤害故意补出抢劫故意。又如年龄跨越发生在结果之前，年龄增长不能倒填早期作为。有排险义务且能防止结果却在达龄后不履行的，另按不作为审查，并证明罪过与因果归责。

\section{第三节\quad 法律认识错误（违法性认识可能性）}

\paragraph{教材模型}
固定行为人 \(p\) 和行为 \(b\)。令 \(K_f\) 为认识到罪状所需的实质事实，\(K_l\) 为知道行为被法律禁止，\(P_l\) 为当时有合理机会认识违法性。
\[
\begin{aligned}
\neg K_f &\Longrightarrow \text{相应的故意认识要素缺失},\\
K_f\land\neg K_l\land P_l &\Longrightarrow
   \text{单凭法律误解不排除责任},\\
K_f\land\neg K_l\land\neg P_l &\Longrightarrow
   \text{依不可避免的违法性错误学说审查责任阻却}.
\end{aligned}
\]
第一行只针对故意所需的事实要素，法定过失犯罪另论。第三行是教材的责任学说，不是刑法总则中的普遍明文免责项。行政许可等规范性要素须区分行为人不了解真实许可状态，与了解状态却误判法律评价。

\paragraph{推导}
甲明知所持物品具有枪支的关键物理特征，明知自己未获任何许可，只误以为私人收藏在法律上享有豁免。题设确认枪支鉴定、适用规定及其他构成要件。甲已认识构成所需的事实，故争点转为 \(P_l\)。如果正常渠道可查明该禁令，仅凭朋友误导不能证明 \(\neg P_l\)。相反，甲误拿与自己行李完全相同的箱子并真信箱内物品归己，错的是归属事实，先审 \(K_f\)。

\paragraph{核验步骤}
有权机关的书面答复可能成为判断 \(P_l\) 的证据，但要核对机关权限、提交事实是否完整、答复对象是否与实际行为相同，以及当时是否还有明显相反的可靠信息。专业意见或书面形式本身不能推出不可避免。知道行为不当也不要求行为人能背出罪名与条号。

\section{第四节\quad 期待可能性}

\paragraph{教材模型和法条边界}
令 \(E(p,t)\) 表示结合当时约束，仍可以期待行为人采取合法替代行为。依学理，若有强制且合理合规路径已被实质封闭，即使已满足形式违法，也应独立评价可谴责性。须先检查正当防卫、紧急避险等法定事由，再检验是否适用第28条胁从犯等明文规则。期待可能性分析不能替代这些法定要件，更不能直接推导普遍的亲属免责。

\paragraph{推导}
甲用眼前紧迫的严重暴力威胁逼乙参与犯罪。已证乙仍有一定选择空间但受强制、实际参与并符合共同犯罪其余条件。由此可评价其受胁迫参加，依法进入第28条的处罚分支，根据情节减轻或者免除。若新事实是甲抓住乙的手强行操作，乙对该动作全无控制，先问是否存在可归属于乙的行为，不能直接认定为胁从犯。若只是一般难堪或人情压力且可安全拒绝，则不足以从这些事实推出 \(\neg E\)。

\paragraph{反例}
父亲明知儿子犯罪，主动为帮助逃匿而提供藏身处，亲属身份并不在《刑法》第310条或两高2021年窝藏、包庇解释中形成一律免责。先分清拒绝作证、积极藏匿、事前通谋等不同事实，再审查相应条文、解释与个案从宽。把道德上的情有可原直接代换成法律上的阻却责任，会漏掉犯罪构成与量刑的区分。

\section*{怎么记与闭卷自测}
\textbf{年龄先分段、能力看当时、错误先问事实、期待找替代。}
年龄在纸上画 \(12,14,16,18,75\) 五道线。前三区间用来决定是否通过年龄门，18与75主要提醒量刑规则。责任能力用两列记录 \(t_b,t_r\)，只有行为列对应能力。认识错误先圈出错认的对象，再问能否认识法律。最后写出一个具体合规方案，用它检验期待可能性。
\begin{enumerate}
\item 13周岁人故意伤害致重伤，未造成严重残疾，也未死亡。若题称情节极恶劣，年龄门能通过吗？\par
答：不能。第17条第3款的结果分支 \(D\lor K\) 不成立，情节不能补替结果门槛。
\item 行为时具有责任能力，结果发生时已丧失能力，是否直接排除责任？\par
答：不能。能力与行为同时审查，结果另证是否由有责任能力时的行为造成。
\item 明知是他人财物却误信刑法允许，与真诚误以为是自己财物，分别先审什么？\par
答：前者在其他要素齐备时审法律认识可能性，后者先审构成事实的认识。
\end{enumerate}
